% Part: propositional-logic

\documentclass[../../include/open-logic-part]{subfiles}

\begin{document}

\olpart{pl}{بياني منطق}

\begin{editorial}
  په دې برخه کښې د کلاسيکي بياني منطق مواد دي۔ لومړے څپرکے
  نسبتاً بنسټيز دے او يوازې تعريفونه او پايلې راغونډوي؛ ډېر ثبوتونه
  نۀ دي بشپړ شوي، بلکې د تمرينونو په توګه پرېښودل شوي دي۔ د ثبوت د
  نظامونو او د بشپړتيا د قضيې مواد د لومړۍ درجې منطق له برخې څخه
  راوړل شوي، چې د لومړۍ درجې منطق ټېګ پکې ناسم ټاکل شوے دے۔ په دې
  توګه د
  محمولونو، اصطلاحاتو او سورونو اړوند هر څه ترې وځي، او د
  !!{structure}s~$\Struct{M}$ پر ځاے د
  !!{valuation}s~$\pAssign{v}$ خبره راځي۔

  پلان دا دے چې دا برخه لا په تفصيل پراخه شي، او نورې موضوعګانې او
  پايلې، لکه د رښتياوالي تابعوي بشپړتيا، هم ورزياتې شي۔
\end{editorial}

\olimport[syntax-and-semantics]{syntax-and-semantics}

\tagfalse{FOL}

\olimport[../first-order-logic/proof-systems]{proof-systems}

\iftag{prfSC}{%
  \olimport[../first-order-logic/sequent-calculus]{sequent-calculus}
}{}

\iftag{prfND}{%
  \olimport[../first-order-logic/natural-deduction]{natural-deduction}
}{}

\iftag{prfTab}{%
  \olimport[../first-order-logic/tableaux]{tableaux}
}{}

\iftag{prfAX}{%
  \olimport[../first-order-logic/axiomatic-deduction]{axiomatic-deduction}
}{}

\olimport[../first-order-logic/completeness]{completeness}

\tagtrue{FOL}

\OLEndPartHook

\end{document}
